\documentclass{beamer}
\usetheme{Madrid}

\usepackage[utf8]{inputenc} % solo si usas pdflatex
\usepackage[T1]{fontenc}    % solo si usas pdflatex

\usepackage{amsmath,amssymb}
\usepackage{microtype}
\usepackage{cancel}         % si lo usas
\usepackage{fontawesome}    % si lo usas
% \usepackage{xypic}        % solo si lo usas
% \usepackage{multicol}     % normalmente no; usar columns de beamer
% \usepackage{amsthm}       % solo si defines teoremas propios
% \usepackage{scrextend}    % solo si lo usas

\usepackage{graphicx}
\graphicspath{{images/}}

\hypersetup{pdfborder={0 0 0}}


\title{Presentación resumen de Java}
\author{Rodrigo Blanco Betes. Alfa Formación}
\date{Enero del 2026}

\begin{document}

\begin{frame}
\titlepage
\end{frame}

\begin{frame}
\tableofcontents
\end{frame}

%-------------------------------------------------


\section{Introducción a Java}
\begin{frame}{Introducción}
Java es un lenguaje de programación de propósito general ampliamente utilizado en:
\begin{itemize}
    \item Aplicaciones de escritorio
    \item Aplicaciones web
    \item Sistemas empresariales
    \item Dispositivos móviles
\end{itemize}

Se caracteriza principalmente por su:
\begin{itemize}
    \item Portabilidad
    \item Seguridad
    \item Robustez
\end{itemize}
\end{frame}

\begin{frame}{Historia de Java}
Java fue desarrollado a principios de la década de 1990 por un equipo liderado por James Gosling en Sun Microsystems.

\begin{itemize}
    \item Proyecto original: \textit{Green Project}
    \item Nombre inicial: \textit{Oak}
    \item Cambio de nombre a \textit{Java} en 1995
    \item Lema: \textit{``Write Once, Run Anywhere''}
\end{itemize}

En 2010, Sun Microsystems fue adquirida por Oracle Corporation, que actualmente mantiene y evoluciona el lenguaje.
\end{frame}

\begin{frame}{Principales Características de Java}
\begin{itemize}
    \item \textbf{Portabilidad}: Compilación a \textit{bytecode} ejecutable en cualquier sistema con JVM.
    \item \textbf{Orientación a objetos}: Facilita la reutilización, mantenimiento y organización del código.
    \item \textbf{Seguridad}: Verificación de bytecode, gestión automática de memoria y modelo de permisos.
    \item \textbf{Gestión automática de memoria}: Uso de recolector de basura.
    \item \textbf{Multihilo}: Soporte para programación concurrente mediante hilos.
\end{itemize}
\end{frame}


\begin{frame}{Conclusión}
Java es un lenguaje maduro y ampliamente utilizado que ha demostrado su utilidad a lo largo del tiempo.

\begin{itemize}
    \item Funciona en múltiples plataformas
    \item Es seguro y robusto
    \item Sigue siendo una opción clave en el desarrollo de software
\end{itemize}
\end{frame}

\section{Tipos primitivos y objetos. API de Java}

\begin{frame}{Tipos primitivos y objetos. API de Java}
  

El lenguaje de programación Java ofrece un conjunto sólido de herramientas fundamentales para el desarrollo de aplicaciones. Entre ellas se encuentran los tipos primitivos, los objetos proporcionados por su API estándar y los instrumentos de control de flujo, que permiten definir la lógica y el comportamiento de los programas.

Los tipos primitivos son los tipos de datos básicos que Java proporciona para almacenar valores simples. Estos \text{no son objetos} y ocupan un espacio fijo en memoria.

\end{frame}


\begin{frame}{Tipos primitivos}

Los tipos primitivos son:
\begin{itemize}
    \item \textbf{byte}: almacena números enteros pequeños.
    \item \textbf{short}: almacena números enteros de tamaño reducido.
    \item \textbf{int}: tipo entero más utilizado.
    \item \textbf{long}: almacena números enteros grandes.
    \item \textbf{float}: almacena números decimales de precisión simple.
    \item \textbf{double}: almacena números decimales de doble precisión.
    \item \textbf{char}: representa un carácter Unicode.
    \item \textbf{boolean}: almacena valores lógicos (\texttt{true} o \texttt{false}).
\end{itemize}

El uso de tipos primitivos permite una mayor eficiencia en el uso de memoria y un mejor rendimiento del programa.


\end{frame}

\begin{frame}{Introducción a la API de Java}

Java incluye una extensa biblioteca estándar conocida como la API de Java, la cual proporciona clases y objetos reutilizables que facilitan el desarrollo de software.

Algunos ejemplos importantes de objetos de la API son:
\begin{itemize}
    \item \textbf{String}: permite manipular cadenas de texto.
    \item \textbf{Scanner}: se utiliza para la entrada de datos desde el teclado.
    \item \textbf{Math}: ofrece métodos para realizar operaciones matemáticas.
    \item \textbf{ArrayList}: permite almacenar colecciones dinámicas de datos.
    \item \textbf{Date}: se utiliza para manejar fechas y tiempos.
\end{itemize}

Estos objetos permiten al programador evitar la implementación desde cero de funcionalidades comunes, aumentando la productividad y reduciendo errores.

\end{frame}

%Estructuras de control de flujo
\section{Estructuras de Control de Flujo}
\begin{frame}{Estructuras de Control de Flujo}
Los estructuras de control de flujo permiten decidir el orden de ejecución de las instrucciones dentro de un programa Java.


Por ejemplo, las estructuras condicionales permiten ejecutar bloques de código según se cumpla o no una condición.
\begin{itemize}
    \item \textbf{if}
    \item \textbf{if-else}
    \item \textbf{switch}
\end{itemize}

Las estructuras repetitivas permiten ejecutar un bloque de código varias veces.
\begin{itemize}
    \item \textbf{for}
    \item \textbf{while}
    \item \textbf{do-while}
\end{itemize}
\end{frame}

\begin{frame}
Estas instrucciones modifican el flujo normal del programa.
\begin{itemize}
    \item \textbf{break}: finaliza un bucle o un \texttt{switch}.
    \item \textbf{continue}: salta a la siguiente iteración de un bucle.
    \item \textbf{return}: devuelve un valor y finaliza la ejecución de un método.
\end{itemize}
\end{frame}

%Arrays
\section{Arrays en Java}
\begin{frame}{Introducción a los arrays}
Un array es una estructura de datos que permite almacenar múltiples valores del mismo tipo bajo un mismo nombre.

Características principales:
\begin{itemize}
    \item Todos los elementos son del mismo tipo
    \item Se accede mediante índices que comienzan en cero
    \item El tamaño es fijo
\end{itemize}

\end{frame}


\begin{frame}[fragile]{Uso de Arrays}

Ejemplos de declaración e inicialización:
    \begin{verbatim}
    int[] numeros = new int[5];

    int[] numeros = {1, 2, 3, 4, 5};
    \end{verbatim}

\end{frame}


\begin{frame}[fragile]{Uso de Arrays}
Acceso a los elementos mediante su índice:
\begin{verbatim}
int valor = numeros[0];
\end{verbatim}

Recorrido de un array con un bucle \texttt{for}:
\begin{verbatim}
for (int i = 0; i < numeros.length; i++) {
    System.out.println(numeros[i]);
}
\end{verbatim}

\begin{verbatim}
for (int numero: numeros) {
    System.out.println(numeros[i]);
}
\end{verbatim}

\end{frame}


\begin{frame}[fragile]{Ventajas y limitaciones de los arrays}

Ventajas y limitaciones:
\begin{itemize}
    \item Acceso rápido y eficiente a los datos
    \item Uso óptimo de la memoria
    \item Tamaño fijo que no puede modificarse
\end{itemize}
\end{frame}


%JOptionPane
\begin{frame}[fragile]{La Clase JOptionPane}
La clase \texttt{JOptionPane} pertenece al paquete \texttt{javax.swing} y permite crear ventanas de diálogo gráficas de forma sencilla, sin interfaces complejas.

Ejemplos de uso:
\begin{itemize}
    \item Mostrar mensajes
    \item Solicitar datos al usuario
    \item Confirmar acciones
\end{itemize}

\end{frame}


\begin{frame}[fragile]{Uso de JOptionPane}

\begin{verbatim}
JOptionPane.showMessageDialog(null, "Hola mundo");

String nombre =
JOptionPane.showInputDialog("Introduce tu nombre");

int opcion =
JOptionPane.showConfirmDialog(null, "¿Deseas continuar?");
\end{verbatim}

Por su simplicidad, es muy recomendable para programas educativos e interacciones sencillas con el usuario.
\end{frame}

\section{Programación Orientada a Objetos en Java}
\begin{frame}{Programación Orientada a Objetos en Java}
Java implementa la Programación Orientada a Objetos (POO) mediante el uso de clases y objetos, permitiendo modelar problemas complejos de forma estructurada.

Los cuatro pilares fundamentales de la POO son:
\begin{itemize}
    \item Encapsulación
    \item Abstracción
    \item Herencia
    \item Polimorfismo
\end{itemize}
\end{frame}

\begin{frame}{Pilares de la POO}
\begin{itemize}
    \item \textbf{Encapsulación}: oculta los detalles internos de una clase y expone solo lo necesario mediante modificadores de acceso.
    \item \textbf{Abstracción}: representa las características esenciales de un objeto, ignorando detalles irrelevantes.
    \item \textbf{Herencia}: permite que una clase herede atributos y métodos de otra.
    \item \textbf{Polimorfismo}: un mismo método puede comportarse de forma distinta según el objeto que lo invoque.
\end{itemize}
\end{frame}

\section{Clases y Objetos en Java}
\begin{frame}{Clases y Objetos}
Una clase es una plantilla que define las propiedades y comportamientos de un objeto.

Un objeto es una instancia concreta de una clase creada a partir de dicha plantilla.
\end{frame}

\begin{frame}[fragile]{Estructura Básica de una Clase}
Una clase en Java se define utilizando la palabra reservada \texttt{class}.

\begin{verbatim}
public class Persona {
    String nombre;
    int edad;
}
\end{verbatim}
\end{frame}

\begin{frame}[fragile]{Creación de Objetos}
Para crear un objeto a partir de una clase se utiliza el operador \texttt{new}.

\begin{verbatim}
Persona p1 = new Persona();
\end{verbatim}
\end{frame}

\section{Constructores}
\begin{frame}[fragile]{Constructores}
Los constructores son métodos especiales que se ejecutan al crear un objeto.

\begin{itemize}
    \item Tienen el mismo nombre que la clase
    \item Se utilizan para inicializar los atributos
\end{itemize}

\begin{verbatim}
public Persona(String nombre, int edad) {
    this.nombre = nombre;
    this.edad = edad;
}
\end{verbatim}
\end{frame}

\section{Métodos}
\begin{frame}[fragile]{Métodos}
Los métodos definen el comportamiento de los objetos y permiten realizar acciones o devolver información.

Ejemplo de método en Java:
\begin{verbatim}
public void mostrarDatos() {
    System.out.println(nombre + " - " + edad);
}
\end{verbatim}
\end{frame}

\section{Uso de Modificadores de Acceso}
\begin{frame}{Modificadores de Acceso}
Java utiliza modificadores de acceso para controlar la visibilidad de atributos y métodos:

\begin{itemize}
    \item \texttt{public}: accesible desde cualquier clase
    \item \texttt{private}: accesible solo desde la misma clase
    \item \texttt{protected}: accesible desde clases del mismo paquete o subclases
\end{itemize}
\end{frame}


\end{document}
